\section{Conclusion}
\label{sec:conclusion}

We introduced Solver Trace Ablation (STA), a protocol that derives the Intensional Structure Ratio (ISR) through leave-one-out component necessity testing on solver inputs.
Our experiments across 6 reward configurations, 3 seeds, 2 formal languages (FOL verified by Prover9, Prolog verified by SWI-Prolog), and up to 1{,}000 training steps establish three findings.

First, ISR is an effective diagnostic for structural reward hacking that prove rate alone cannot detect.
Under execution-only GRPO, prove rate rises from 32\% to 96\% while ISR collapses from 0.64 to 0.003, with 94.8\% of proved outputs at 1{,}000 steps containing zero necessary components.
All three seeds exhibit 87--100\% conclusion embedding, and a supervised fine-tuning baseline (92.9\% genuine, near-zero hacking) confirms that these pathologies are RL-reward-specific.
The joint signal of prove rate, ISR, and component count cleanly separates three qualitatively different failure modes (redundant embedding, repetition collapse, tautological collapse) that prove rate conflates into apparent success.

Second, ISR fails as a training signal due to meta-level Goodhart's Law.
At low dose ($\beta{=}0.1$), ISR reward induces ISR-specific gaming: 64.5\% of proved outputs across 3 seeds exhibit single-premise conclusion duplication (H7), compared to 0\% across all random-reward seeds.
The apparent structural advantage of ISR reward (raw 2.4$\times$ genuine ISR) is entirely a premise minimization artifact---controlling for premise count yields 1.06$\times$, indistinguishable from the random baseline.
At high dose ($\beta{=}0.3$), classic Goodhart collapse drives prove rate to 18\% versus 48\% for the matched random control.
Conclusion embedding itself is GRPO-universal (68--100\% across conditions); ISR reward restructures hacking from multi-premise (H1) to single-premise (H7) forms without reducing its prevalence.

Third, both the diagnostic and the Goodhart failure transfer across formal languages.
In Prolog, execution-only training again produces near-zero ISR (mean 0.04, 3 seeds), and ISR-as-reward produces universal Goodhart: all 3 seeds converge to gaming with 0\% genuine output, each via a distinct pathway---tautological duplication, degenerate uppercase syntax, and query-as-fact gaming.
The stronger Prolog result (0\% genuine vs.\ ${\sim}$16\% in FOL) suggests that simpler formal semantics enable more complete metric exploitation.

\paragraph{Limitations.}
Our experiments use a single base model (Qwen3-4B); whether larger models exhibit qualitatively different hacking dynamics remains open.
The evaluation covers two formal languages (FOL and Prolog) and two datasets (FOLIO-derived and ProofWriter); broader logic types (e.g., dependent types, temporal logic) may present different reward hacking landscapes.
ISR measures syntactic component necessity, not semantic faithfulness---downstream NLI accuracy shows no correlation with ISR ($\rho{=}{-}0.049$, n.s.)---and the LOO formulation carries an inherent premise-count bias that requires normalization.
We evaluate with 3 seeds; while the cross-seed divergence is itself a finding, larger seed counts would strengthen statistical claims.
The text-match hacking classifier has blind spots exposed by Prolog (query-as-fact gaming bypassed strict matching), necessitating functional cross-validation via $n_{\text{necessary}}$ distributions.

\paragraph{Future directions.}
Combining ISR with semantic verification (e.g., NLI-based faithfulness checks) could address the syntactic-only limitation and yield diagnostics sensitive to both structural and semantic quality.
Adaptive reward weighting that uses ISR as a gating signal rather than a continuous bonus may avoid the meta-Goodhart trap.
Extending STA to richer logic types (SMT solvers, interactive theorem provers) and scaling to larger models would test the generality of the hacking dynamics we document.
Finally, the recursive Goodhart phenomenon---where diagnostics designed to detect gaming become gaming targets when optimized---warrants investigation as a general failure mode of proxy-metric-as-reward approaches beyond the neurosymbolic setting.
